% !TeX root = ../main.tex
\chapter{Projekt systemu}
\label{chap:design}

% FIRST DRAFT. Based on docs/architecture.md, the ADRs and the PAI design
% documents (architecture, logic, data structures, user interface).

Projekt systemu wyrasta z jednej obserwacji: opinia powinna przejść od źródła do wykresu drogą, którą da się powtórzyć.
Każdy etap tej drogi zostawia trwały ślad, z którego można odtworzyć wszystkie następne.
W tym rozdziale opisujemy architekturę systemu i uzasadnienie każdego jej elementu, logikę odtworzenia, model danych w hurtowni, projekt panelu oraz sposób, w jaki system chroni prywatność autorów opinii.

\section{Architektura}
\label{sec:design-architecture}

Drogę pojedynczej opinii przez system przedstawia rysunek~\ref{fig:architecture}.
Opinię pobiera kolektor i zapisuje ją w archiwum surowych danych.
Z archiwum opinie trafiają do hurtowni, skąd wydawca odtworzeń publikuje je w kolejce komunikatów w porządku chronologicznym.
Potok przetwarzania strumieniowego odczytuje komunikaty, wykrywa język, zleca klasyfikację wydźwięku modelowi na punkcie końcowym Agent Platform i dopisuje wynik do tabeli lądowania.
Po zakończeniu odtworzenia wyniki są scalane z docelową tabelą zdarzeń, na której opierają się widoki raportowe i panel.

\begin{figure}[hbp]
  \centering
  \begin{tikzpicture}[
      font=\footnotesize,
      box/.style={draw, rounded corners=2pt, align=center, minimum height=10mm, text width=23mm, fill=gray!8},
      store/.style={box, fill=blue!8},
      ext/.style={box, dashed, fill=white},
      arr/.style={-{Latex[length=2mm]}, thick},
      lab/.style={font=\scriptsize, fill=white, inner sep=1pt}]
    \node[ext, minimum height=8mm] (yt) at (0,0.65) {YouTube\\Data API};
    \node[ext, minimum height=8mm] (st) at (0,-0.65) {serwis recenzji\\Steam};
    \node[box] (col) at (3.8,0) {kolektory\\(Cloud Run Jobs)};
    \node[store] (gcs) at (7.6,0) {archiwum\\surowych danych\\(Cloud Storage)};
    \node[store] (stg) at (11.4,0) {BigQuery:\\\texttt{raw\_landing},\\\texttt{raw\_staging}};
    \node[box] (pub) at (11.4,-2.7) {wydawca\\odtworzeń\\(Cloud Run Job)};
    \node[box] (ps) at (7.6,-2.7) {temat Pub/Sub\\(klucz \enquote{global})};
    \node[box] (df) at (3.8,-2.7) {potok\\Apache Beam\\(Dataflow)};
    \node[box] (vx) at (0,-2.7) {punkt końcowy\\Agent Platform\\(XLM-R)};
    \node[box] (dlq) at (0,-5.4) {temat błędnych\\komunikatów};
    \node[store] (ev) at (3.8,-5.4) {BigQuery:\\\texttt{events\_landing},\\\texttt{events}};
    \node[store] (rep) at (7.6,-5.4) {widoki\\raportowe\\(Dataform)};
    \node[ext] (ls) at (11.4,-5.4) {panel\\(Data Studio)};
    \draw[arr] (yt) -- (col);
    \draw[arr] (st) -- (col);
    \draw[arr] (col) -- node[lab, above=1pt] {JSONL.gz} (gcs);
    \draw[arr] (gcs) -- node[lab, above=1pt] {ładowanie} (stg);
    \draw[arr] (stg) -- node[lab, right=1pt, align=left] {ORDER BY\\created\_utc} (pub);
    \draw[arr] (pub) -- node[lab, above=1pt] {publikacja} (ps);
    \draw[arr] (ps) -- (df);
    \draw[arr, {Latex[length=2mm]}-{Latex[length=2mm]}] (df) -- node[lab, above=1pt] {REST} (vx);
    \draw[arr] (df) -- node[lab, right=1pt] {dopisanie} (ev);
    \draw[arr] (df) -- (dlq);
    \draw[arr] (ev) -- node[lab, above=1pt] {MERGE} (rep);
    \draw[arr] (rep) -- (ls);
  \end{tikzpicture}
  \caption{Droga opinii przez system: od serwisów źródłowych do panelu analitycznego}
  \label{fig:architecture}
\end{figure}

Każdy z tych etapów istnieje z konkretnego powodu.
Poniżej opisujemy te powody, bo bez nich architektura mogłaby wyglądać na nadmiernie rozbudowaną.

\subsection{Dlaczego archiwum obok hurtowni}
\label{sec:design-archive}

Hurtownia BigQuery jest analitycznym źródłem prawdy, ale nie jedynym trwałym zapisem danych.
Archiwum w Cloud Storage przechowuje opinie dokładnie w takiej postaci, w jakiej je pobrano, jako skompresowane pliki z jednym obiektem \ac{json} w każdym wierszu.
Pliki trafiają pod ścieżkę złożoną ze źródła, wydarzenia, daty i godziny pobrania oraz unikalnego identyfikatora paczki, więc dwa przebiegi kolektora nigdy nie nadpisują swoich wyników.
Jeśli zmieni się model albo schemat danych, hurtownię odbudowujemy z archiwum, a nie z internetu.
Interfejsy serwisów mogą do tego czasu zniknąć, zmienić się albo wyczerpać limit, a komentarze mogą zostać usunięte przez autorów.

\subsection{Dlaczego wydawca między hurtownią a kolejką}
\label{sec:design-publisher}

Globalna kolejność chronologiczna opinii ze wszystkich źródeł jest wymaganiem, a Pub/Sub sam jej nie zapewnia.
Zapewnia ją natomiast BigQuery: zapytanie z klauzulą \code{ORDER BY created_utc, id} sortuje cały zbiór, a identyfikator rozstrzyga remisy, gdy dwie opinie mają ten sam czas z dokładnością do sekundy.
Identyfikator zaczyna się od nazwy źródła, więc przy remisie recenzje Steam zawsze poprzedzają komentarze YouTube i kolejność jest za każdym razem taka sama.
Wydawca odczytuje posortowany wynik i publikuje rekordy pod jednym, stałym kluczem porządkującym \code{global}, więc Pub/Sub dostarcza je w tej samej kolejności.
Czas opublikowania opinii przekazuje jako atrybut komunikatu, z którego potok odczytuje czas zdarzenia.

Wydawca nie wysyła rekordów tak szybko, jak potrafi.
Odtwarza rzeczywiste odstępy między opiniami, podzielone przez współczynnik przyspieszenia, a odstępy dłuższe niż zadany próg skraca do tego progu.
Dzięki temu skupiska opinii wokół wydarzeń zachowują swój kształt, a miesiące ciszy między wydarzeniami znikają.
Strumień wygląda jak żywy, tylko szybszy, a przebieg w czasie jest wciąż ten sam.

\subsection{Dlaczego tabele lądowania przed tabelami docelowymi}
\label{sec:design-landing}

Pub/Sub i Dataflow gwarantują dostarczenie komunikatu co najmniej raz, a nie dokładnie raz.
Zamiast walczyć z duplikatami w potoku, przyjęliśmy, że potok zapisuje wyłącznie do tabeli lądowania, do której tylko dopisuje.
Dopiero instrukcja \code{MERGE} przenosi wyniki do tabeli docelowej, zostawiając jeden wiersz na identyfikator opinii, ten z najpóźniejszym czasem przetworzenia.
Tak samo działa ładowanie archiwum: pliki trafiają najpierw do tabeli ładowania, a stamtąd, bez duplikatów, do tabeli pośredniej.
Dzięki temu odtworzenie można powtarzać dowolnie wiele razy, a tabela docelowa nie rośnie.

\subsection{Dlaczego model na zewnątrz potoku}
\label{sec:design-remote-model}

Potok nie zawiera wag modelu.
Wywołuje klasyfikator przez wspólny interfejs, za którym może stać deterministyczna zaślepka albo klient punktu końcowego Agent Platform.
Taki podział ma dwie przyczyny.
Maszyny robocze potoku działają bez publicznych adresów IP i łączą się z usługami Google wyłącznie przez prywatny dostęp w sieci \ac{vpc}, więc nie mogą pobrać wag z internetu przy starcie.
Gdybyśmy z kolei wbudowali wagi w obraz potoku, każda zmiana modelu wymagałaby przebudowania obrazu, a obraz rozrósłby się o kilka gigabajtów.
Punkt końcowy pozwala zmienić model bez dotykania potoku i bez ponownego budowania czegokolwiek.

\subsection{Dlaczego panel czyta widoki, a nie tabele}
\label{sec:design-views}

Panel ma jedno źródło danych dla każdego widoku w zbiorze raportowym.
Widoki zawierają gotowe miary, zdefiniowane w języku \ac{sql} i sprawdzane testami przy każdym wdrożeniu, a panel tylko sumuje liczniki i dzieli je raz.
Dzięki temu liczba na wykresie zawsze zgadza się z wynikiem zapytania.
Widoki są też granicą prywatności: nie zawierają pseudonimu autora ani treści opinii, a osoba oglądająca panel nie ma dostępu do tabeli zdarzeń.

\subsection{Wdrożenie i komponenty}
\label{sec:design-deployment}

Rysunek~\ref{fig:deployment} pokazuje rozmieszczenie komponentów w chmurze.
Wszystkie działają w projekcie Google Cloud w regionie \code{europe-central2}, a użytkownik potrzebuje tylko przeglądarki.
Diagram przedstawia stan z października 2026~r., gdy panel działał jeszcze jako tymczasowa aplikacja w Cloud Run.
\todo{Zaktualizować diagram wdrożenia: usunąć panel tymczasowy, dodać serwer MLflow (Cloud Run i Cloud SQL).}

\begin{figure}[hbp]
  \centering
  \includegraphics[width=0.62\textwidth]{deployment}
  \caption{Diagram wdrożenia systemu}
  \label{fig:deployment}
\end{figure}

Z punktu widzenia oprogramowania system ma trzy warstwy.
Warstwa prezentacji obejmuje panel i interfejs operacyjny zespołu, czyli konsolę Google Cloud, polecenia \code{gcloud} i powiadomienia e-mail.
Warstwa usług obejmuje kolektory, wydawcę odtworzeń, potok, moduł trenowania i orkiestrację odtworzenia.
Warstwa danych obejmuje archiwum, oba zbiory BigQuery, tematy Pub/Sub, repozytorium modeli oraz konfigurację i sekrety.
Usługi nie wywołują się nawzajem bezpośrednio, tylko wymieniają dane przez warstwę danych.
Jedynym wyjątkiem jest wywołanie klasyfikatora przez potok.
Podział na podsystemy i przepływ danych między nimi przedstawia rysunek~\ref{fig:components}.

\begin{figure}[hbp]
  \centering
  \includegraphics[width=0.9\textwidth]{components-subsystems}
  \caption{Diagram komponentów w podziale na podsystemy}
  \label{fig:components}
\end{figure}

\section{Logika odtworzenia}
\label{sec:design-replay}

Odtworzenie jest operacją, wokół której zbudowany jest cały system, więc jego logikę zaprojektowaliśmy szczególnie starannie.
Jedno odtworzenie to jedno uruchomienie wyzwalacza w Cloud Build, które wykonuje sześć kroków: uruchamia potok, publikuje strumień, czeka na komplet wyników, zatrzymuje potok, scala wyniki i wylicza odcisk tabeli.
Kolejność tych kroków i warunek, od którego zależy scalenie, przedstawia rysunek~\ref{fig:sequence-replay}.

\begin{figure}[hbp]
  \centering
  \includegraphics[width=0.82\textwidth]{sequence-replay}
  \caption{Diagram sekwencji uruchomienia odtworzenia}
  \label{fig:sequence-replay}
\end{figure}

Logikę odtworzenia wyznaczają trzy reguły.
Po pierwsze, zatrzymanie potoku czeka na kompletność, a nie na ciszę.
Orkiestrator wylicza z tabeli pośredniej oczekiwaną liczbę rekordów i czeka, aż tyle różnych identyfikatorów pojawi się w tabeli lądowania z czasem przetworzenia późniejszym niż początek bieżącego odtworzenia.
Licznik, który przestał rosnąć, nie jest dowodem zakończenia, bo zatrzymany potok wygląda dokładnie tak samo.
Po drugie, niekompletne odtworzenie nigdy nie jest scalane: jeśli rekordy nie dotrą w zadanym czasie, orkiestrator zatrzymuje potok i kończy się błędem, a tabela zdarzeń pozostaje bez zmian.
Po trzecie, potok jest zatrzymywany zawsze, także wtedy, gdy wcześniejszy krok się nie powiódł, bo zadanie strumieniowe, którego nikt nie zatrzymał, nalicza opłaty bez końca.

Stany, przez które przechodzi odtworzenie, przedstawia rysunek~\ref{fig:state-replay}.
Odtworzenie przechodzi je automatycznie, a błąd na dowolnym etapie przenosi je do stanu przerwania.

\begin{figure}[hbp]
  \centering
  \includegraphics[width=0.42\textwidth]{state-replay}
  \caption{Diagram stanów odtworzenia}
  \label{fig:state-replay}
\end{figure}

Ostatnim krokiem odtworzenia jest odcisk tabeli zdarzeń, czyli suma skrótów wszystkich wierszy z pominięciem kolumn, które z definicji różnią się między przebiegami: wersji modelu i czasu przetworzenia.
Suma nie zależy od kolejności wierszy, więc dwa odtworzenia tych samych danych tym samym klasyfikatorem muszą dać tę samą liczbę.
Odcisk zapisuje się w dzienniku uruchomienia razem z wersją kodu i parametrami, a dziennik ten jest dowodem pochodzenia wyników.
Wyzwalacz odtworzenia uruchamia się wyłącznie ręcznie, bo żadne scalenie zmian w repozytorium nie powinno samo uruchamiać najdroższego zasobu w projekcie.

\section{Model danych}
\label{sec:design-data-model}

Hurtownia ma dwa zbiory danych: analityczny, z tabelami, i raportowy, z widokami.
Zbiór analityczny przedstawia rysunek~\ref{fig:db-analytics}.
Zawiera cztery tabele o dwóch schematach.
Tabele \code{raw_landing} i \code{raw_staging} przechowują rekordy opinii w formacie z tabeli~\ref{tab:record-schema}.
Tabele \code{events_landing} i \code{events} przechowują rekordy wzbogacone o wykryty język, etykietę wydźwięku, pewność modelu, wersję modelu i czas przetworzenia.

\begin{figure}[hbp]
  \centering
  \includegraphics[width=0.6\textwidth]{db-analytics}
  \caption{Struktura zbioru analitycznego w BigQuery}
  \label{fig:db-analytics}
\end{figure}

Świadomie zrezygnowaliśmy z normalizacji.
Tabela \code{events} jest jedną zdenormalizowaną tabelą, w której każda opinia ma wszystkie potrzebne kolumny, więc żaden widok raportowy nie musi wykonywać złączeń.
Tabele pośrednia i docelowa są partycjonowane dziennie według czasu publikacji opinii i klastrowane według źródła i wydarzenia.
Przy obecnej liczbie wierszy to przesada, ale zapytania z filtrem czasowym czytają dzięki temu tylko potrzebne partycje, a ten sam schemat skaluje się do znacznie większych zbiorów.

Zbiór raportowy zawiera osiem widoków, zestawionych w tabeli~\ref{tab:reporting-views}.
Widoki są autoryzowane na zbiorze analitycznym jako całości, więc mogą czytać tabelę \code{events}, do której oglądający panel nie ma dostępu.

\begin{table}[hbp]
  \caption{Widoki zbioru raportowego}
  \label{tab:reporting-views}
  \centering
  \small
  \begin{tabular}{|l|p{0.58\textwidth}|}
    \hline
    Widok                                    & Zawartość                                                                                       \\
    \hline
    \texttt{v\_events}                       & jeden wiersz na opinię, bez pseudonimu autora i treści; podstawa pozostałych widoków             \\
    \hline
    \texttt{v\_sentiment\_daily}             & liczby opinii w klasach wydźwięku dla dnia, źródła i wydarzenia                                  \\
    \hline
    \texttt{v\_sentiment\_by\_period}        & to samo w ujęciu dziennym, tygodniowym i miesięcznym, z zerami dla okresów bez danych            \\
    \hline
    \texttt{v\_sentiment\_by\_event}         & zestawienie dla wydarzeń: liczba opinii, udziały klas, wskaźnik netto, zakres dat                \\
    \hline
    \texttt{v\_sentiment\_by\_event\_source} & to samo w podziale na źródła                                                                     \\
    \hline
    \texttt{v\_event\_lift}                  & efekt wydarzenia: zmiana wskaźnika netto i liczby opinii na dzień względem tła z tego samego źródła \\
    \hline
    \texttt{v\_source\_divergence}           & różnica wskaźnika netto między najbardziej i najmniej pozytywnym źródłem dla wydarzenia          \\
    \hline
    \texttt{v\_source\_coverage}             & liczba opinii i udział pozytywnych w podziale na źródła i języki                                 \\
    \hline
  \end{tabular}
\end{table}

\section{Projekt interfejsu użytkownika}
\label{sec:design-ui}

Panel składa się z trzech stron, odpowiadających trzem przypadkom użycia analityka.
Strona przeglądu pokazuje wskaźniki podsumowujące, czyli liczbę opinii, wskaźnik netto oraz udziały opinii pozytywnych i negatywnych, a pod nimi wykres wskaźnika netto w czasie dla źródeł wybranych w filtrze i wykres liczby opinii.
Strona wydarzeń zestawia wskaźnik netto i podział na klasy dla każdego wydarzenia i źródła, różnice między źródłami oraz efekt wydarzenia względem tła.
Strona źródeł i języków pokazuje liczbę opinii i udział opinii pozytywnych w podziale na języki.

Przy projektowaniu stron kierowaliśmy się trzema zasadami.
Przy każdym wyniku widać liczbę opinii, z której go policzono, bo wskaźnik netto z dziesięciu komentarzy znaczy co innego niż ten sam wskaźnik z dziesięciu tysięcy.
Na wykresach w czasie pokazujemy tylko okresy z co najmniej dwudziestoma opiniami, a daty okien wydarzeń podaje tabela pod wykresami.
Kolory źródeł i klas wydźwięku są stałe na wszystkich stronach.

Strony najpierw zaimplementowaliśmy jako aplikację w bibliotece Streamlit, a następnie odtworzyliśmy w Data Studio, z tymi samymi widokami i tymi samymi liczbami.
Rysunek~\ref{fig:ui-overview} przedstawia górną część strony przeglądu w raporcie Data Studio: filtry, cztery wskaźniki podsumowujące i wykres wskaźnika netto w czasie.
Pozostałe części raportu pokazujemy w podrozdziale~\ref{sec:guide-dashboard}.

\begin{figure}[hbp]
  \centering
  \includegraphics[width=\textwidth]{ui-overview-top}
  \caption{Górna część strony \enquote{Przegląd} w raporcie Data Studio}
  \label{fig:ui-overview}
\end{figure}

\section{Prywatność i bezpieczeństwo}
\label{sec:design-privacy}

Identyfikatory autorów opinii, takie jak nazwa kanału YouTube czy numer konta Steam, są danymi osobowymi w rozumieniu \ac{rodo}~\cite{rodo2016}.
Dlatego granicę prywatności postawiliśmy możliwie wcześnie, w samym kolektorze, zanim rekord zostanie gdziekolwiek zapisany.
Kolektor zastępuje identyfikator autora skrótem SHA-256 z identyfikatora i tajnej soli, przyciętym do 16 znaków szesnastkowych.
Ten sam autor ma zawsze ten sam pseudonim, co pozwala liczyć unikalnych autorów, ale bez soli nie da się z pseudonimu odtworzyć identyfikatora.
Sól jest przechowywana w Secret Manager i nie może być zmieniana, bo zmiana zerwałaby ciągłość pseudonimów w zbiorze.
W przypadku Steam kolektor odrzuca też nazwę wyświetlaną autora i adres jego profilu.

Tak przetworzone dane są w rozumieniu \ac{rodo} danymi spseudonimizowanymi, a nie anonimowymi, więc nadal podlegają ochronie.
Dlatego druga granica przebiega między zbiorem analitycznym a raportowym.
Widoki raportowe nie zawierają ani pseudonimu, ani treści opinii, a liczba unikalnych autorów jest liczona wewnątrz agregatów.
Test \code{no_personal_columns}, uruchamiany przy każdym wdrożeniu widoków, przerywa wdrożenie, gdy którykolwiek obiekt w zbiorze raportowym ma kolumnę o nazwie \code{author_hash} albo \code{text}.

Pozostałe zasady bezpieczeństwa są typowe dla systemów w chmurze, ale konsekwentnie stosowane.
Każdy komponent ma osobne konto usługi z minimalnym zestawem uprawnień, a żadne konto nie jest współdzielone.
Maszyny robocze potoku i środowisko trenowania nie mają publicznych adresów IP.
Sekrety, takie jak klucz YouTube Data API i sól, istnieją wyłącznie w Secret Manager, nigdy w kodzie ani w dziennikach.
Serwer MLflow jest dostępny tylko przez \ac{iap}, dla osób z listy, a raport w Data Studio wykonuje zapytania z uprawnieniami osoby, która go ogląda, więc samo udostępnienie raportu nie daje dostępu do danych.
Wszystkie dane przetwarzamy i przechowujemy w regionie warszawskim.
Jak te założenia przełożyły się na kod, opisuje następny rozdział.
